\documentclass[11pt]{article}
\usepackage[margin=1in]{geometry}
\usepackage{amsmath,amssymb,amsthm}
\usepackage{mathrsfs}
\usepackage{graphicx}
\usepackage[hidelinks]{hyperref}

\newtheorem{theorem}{Theorem}
\newtheorem{lemma}{Lemma}
\newtheorem{corollary}{Corollary}
\newtheorem{remark}{Remark}

\newcommand{\R}{\mathbb{R}}
\newcommand{\D}{\mathbb{D}}
\newcommand{\g}{\gamma}
\newcommand{\cL}{\mathcal L}

\begin{document}

\begin{center}
{\Large A separable-domain answer for the Lipschitz Malliavin chain rule}\\[4pt]
{\normalsize Historical partial result for arXiv:1204.2946}
\end{center}

\section*{Status}

This is a reconstructed predecessor packet preserved under the history folder
of the full basis-free solution.  It is superseded by the parent packet, which
removes the separability restriction from the operator-factor conclusion.

\section*{Source Question}

Pronk and Veraar prove in \cite[Proposition 3.10]{PV} that if \(E_0\) has a
Schauder basis, \(E_1\) is UMD, \(p\in(1,\infty)\), and
\(\phi:E_0\to E_1\) is Lipschitz, then
\[
  F\in \D^{1,p}(E_0) \quad\Longrightarrow\quad
  \phi(F)\in \D^{1,p}(E_1),
\]
and there is a bounded linear operator
\[
  T_F:\g(H,E_0)\to L^\infty(\Omega;\g(H,E_1))
\]
such that \(D(\phi(F))=T_F(DF)\).  In Remark 3.11 they write that they do not
know whether the assumption that \(E_0\) has a basis can be avoided.

\begin{center}
\includegraphics[width=0.95\linewidth]{figures/open_problem_crop.png}
\end{center}

\section*{Historical Result}

\begin{theorem}\label{thm:membership}
Let \(H\) be a separable Hilbert space, let \(p\in(1,\infty)\), let \(E_0\) be
an arbitrary Banach space, and let \(E_1\) be a UMD Banach space.  If
\(\phi:E_0\to E_1\) is Lipschitz with constant \(L\), then for every
\(F\in\D^{1,p}(E_0)\) one has
\[
  \phi(F)\in \D^{1,p}(E_1).
\]
Moreover
\[
  \|D\phi(F)(\omega)\|_{\g(H,E_1)}
  \leq L\|DF(\omega)\|_{\g(H,E_0)}
  \quad\hbox{for a.e. }\omega\in\Omega .
\]
\end{theorem}

\begin{theorem}\label{thm:separable-factor}
In the setting of Theorem \ref{thm:membership}, assume in addition that
\(E_0\) is separable.  Then there is a bounded linear operator
\[
  T_F:\g(H,E_0)\to L^\infty(\Omega;\g(H,E_1))
\]
with \(\|T_F\|\leq L\) such that
\[
  D(\phi(F))=T_F(DF).
\]
\end{theorem}

\begin{corollary}
The Schauder-basis assumption in \cite[Proposition 3.10]{PV} may be replaced
by separability of \(E_0\).  In particular, no approximation property of
\(E_0\) is needed for the full operator-factor conclusion in the separable
case.
\end{corollary}

\section*{Proof Intuition}

The source proof uses a basis to approximate the identity on \(E_0\) by
uniformly bounded finite-dimensional projections.  This is stronger than what
the chain rule needs.  A Malliavin-Sobolev random variable is already obtained
as a limit of smooth random variables with finite-dimensional ranges.  We can
therefore apply the finite-dimensional Lipschitz chain rule to those
approximants and pass to the limit.

The remaining operator \(T_F\) is a factorization of the already constructed
derivative through \(DF\).  Once the pointwise estimate
\(\|D\phi(F)\|\leq L\|DF\|\) is known, separability of \(\g(H,E_0)\) gives a
measurable norming functional for \(DF\).  This norming functional produces a
rank-one operator field that maps \(DF(\omega)\) to \(D\phi(F)(\omega)\) with
norm at most \(L\).

\section*{A finite-dimensional lemma}

\begin{lemma}\label{lem:finite}
Let \(V\) be a finite-dimensional Banach space, let \(E_1\) be UMD, and let
\(\psi:V\to E_1\) be Lipschitz with constant \(L\).  If
\(G\in\D^{1,p}(V)\), then \(\psi(G)\in\D^{1,p}(E_1)\) and
\[
  \|D\psi(G)(\omega)\|_{\g(H,E_1)}
  \leq L\|DG(\omega)\|_{\g(H,V)}
  \quad\hbox{a.s.}
\]
\end{lemma}

\begin{proof}
Identify \(V\) linearly with \(\R^m\) and convolve \(\psi\) with a standard
smooth mollifier on \(V\).  This gives \(C^\infty\) maps
\(\psi_\varepsilon:V\to E_1\) such that
\[
  \|\psi_\varepsilon(x)-\psi(x)\|\to0\quad (x\in V),
  \qquad
  \|\psi_\varepsilon'(x)\|_{\cL(V,E_1)}\leq L .
\]
The smooth chain rule \cite[Proposition 3.9]{PV} gives
\[
  \psi_\varepsilon(G)\in\D^{1,p}(E_1),\qquad
  D\psi_\varepsilon(G)=\psi_\varepsilon'(G)DG,
\]
and the ideal property of \(\gamma\)-radonifying operators gives the pointwise
bound
\[
  \|D\psi_\varepsilon(G)(\omega)\|_{\g(H,E_1)}
  \leq L\|DG(\omega)\|_{\g(H,V)} .
\]
Also \(\psi_\varepsilon(G)\to\psi(G)\) in \(L^p(\Omega;E_1)\).  Since \(E_1\)
is UMD, \(\g(H,E_1)\) is reflexive, so the compactness lemma
\cite[Lemma 3.8]{PV} implies \(\psi(G)\in\D^{1,p}(E_1)\), after passing to a
weakly convergent derivative subsequence.

Mazur's lemma gives convex combinations of the derivatives
\(D\psi_{\varepsilon_j}(G)\) which converge strongly in
\(L^p(\Omega;\g(H,E_1))\) to \(D\psi(G)\).  The same pointwise bound is
preserved by convex combinations.  Passing to an almost surely convergent
subsequence yields the asserted pointwise estimate.
\end{proof}

\section*{Proof of Theorem \ref{thm:membership}}

Choose smooth random variables \(F_n\in\mathscr S\otimes E_0\) such that
\[
  F_n\to F \quad\hbox{in }\D^{1,p}(E_0).
\]
Each \(F_n\) takes values in a finite-dimensional subspace \(V_n\subseteq E_0\).
Applying Lemma \ref{lem:finite} to \(\phi|_{V_n}:V_n\to E_1\) gives
\[
  \phi(F_n)\in\D^{1,p}(E_1),\qquad
  \|D\phi(F_n)(\omega)\|_{\g(H,E_1)}
  \leq L\|DF_n(\omega)\|_{\g(H,E_0)} .
\]
Since \(\phi\) is Lipschitz,
\[
  \phi(F_n)\to\phi(F) \quad\hbox{in } L^p(\Omega;E_1).
\]
The derivatives \(D\phi(F_n)\) are bounded in
\[
  L^p(\Omega;\g(H,E_1)),
\]
so the compactness lemma again gives
\(\phi(F)\in\D^{1,p}(E_1)\) and, after passing to a subsequence,
\[
  D\phi(F_n)\rightharpoonup D\phi(F)
  \quad\hbox{in }L^p(\Omega;\g(H,E_1)).
\]

It remains to keep the pointwise estimate.  By Mazur's lemma there are convex
combinations
\[
  Z_k=\sum_{n\geq k} a_{k,n}D\phi(F_n)
\]
which converge strongly in \(L^p(\Omega;\g(H,E_1))\) to \(D\phi(F)\).  Set
\[
  M_k=\sum_{n\geq k} a_{k,n}\|DF_n\|_{\g(H,E_0)} .
\]
Since \(DF_n\to DF\) in \(L^p(\Omega;\g(H,E_0))\), we have
\[
  M_k\to \|DF\|_{\g(H,E_0)}
  \quad\hbox{in }L^p(\Omega).
\]
For every \(k\),
\[
  \|Z_k(\omega)\|_{\g(H,E_1)}\leq L M_k(\omega)
  \quad\hbox{a.s.}
\]
Taking an almost surely convergent subsequence of \(Z_k\) and \(M_k\) gives
\[
  \|D\phi(F)(\omega)\|_{\g(H,E_1)}
  \leq L\|DF(\omega)\|_{\g(H,E_0)}
\]
for almost every \(\omega\).  This proves Theorem \ref{thm:membership}.

\section*{Proof of Theorem \ref{thm:separable-factor}}

Put
\[
  G=\g(H,E_0),\qquad Y=\g(H,E_1),\qquad X(\omega)=DF(\omega),\qquad
  Z(\omega)=D\phi(F)(\omega).
\]
If \(E_0\) and \(H\) are separable, then \(G\) is separable because finite-rank
operators with coefficients from countable dense subsets are dense in
\(\g(H,E_0)\).

We use the standard measurable norming-functional selection for separable
Banach spaces: there is a Borel map
\[
  J:G\to B_{G^*}
\]
such that \(Jx(x)=\|x\|\) for all \(x\in G\).  This follows from the
Kuratowski--Ryll-Nardzewski selection theorem applied to the compact-valued map
\(x\mapsto\{x^*\in B_{G^*}:x^*(x)=\|x\|\}\), where \(B_{G^*}\) is equipped with
the weak-star topology.

For \(R\in G\), define
\[
  (T_FR)(\omega)=
  \begin{cases}
    J(X(\omega))(R)\, Z(\omega)/\|X(\omega)\|, & X(\omega)\neq0,\\[3pt]
    0, & X(\omega)=0 .
  \end{cases}
\]
The map \(\omega\mapsto J(X(\omega))(R)\) is measurable for every fixed
\(R\), and \(Z\) is strongly measurable, so \(T_FR\) is strongly measurable as
a \(Y\)-valued random variable.  The pointwise estimate from
Theorem \ref{thm:membership} gives
\[
  \|(T_FR)(\omega)\|_Y\leq L\|R\|_G
\]
almost surely; hence \(T_F\) maps \(G\) boundedly into
\(L^\infty(\Omega;Y)\) and \(\|T_F\|\leq L\).  Linearity in \(R\) is pointwise
linearity of \(J(X(\omega))\).  Finally,
\[
  (T_FX)(\omega)
  =J(X(\omega))(X(\omega))\, Z(\omega)/\|X(\omega)\|
  =Z(\omega)
\]
when \(X(\omega)\neq0\), while if \(X(\omega)=0\) then \(Z(\omega)=0\) by the
same pointwise estimate.  Therefore \(T_F(DF)=D\phi(F)\).

\section*{Historical Limitation}

This predecessor removed the Schauder-basis and bounded approximation property
requirements in the separable version of the source proposition, and proved
the membership theorem for arbitrary \(E_0\).  It remained partial because the
operator-factor conclusion was proved only under separability of \(E_0\).
The parent packet supersedes this limitation.

\begin{thebibliography}{9}
\bibitem{PV}
M. Pronk and M. Veraar,
\emph{Tools for Malliavin calculus in UMD Banach spaces},
Potential Analysis 40 (2014), 307--344; arXiv:1204.2946.
\end{thebibliography}

\end{document}
